\subsection{Análisis de Eficiencia Computacional, Latencias y Huella de Almacenamiento}
\label{subsec:eficiencia_computacional_costes}

Esta subsección analiza la viabilidad computacional y operativa de la arquitectura propuesta, evaluando los desgloses de latencia por etapa, la escalabilidad de almacenamiento vectorial y los compromisos de diseño para su adopción en entornos de producción.

\subsubsection{Desglose Granular de Latencias}
El arnés de auditoría \texttt{RunAuditor} registró los tiempos de ejecución de cada componente del prototipo:
\begin{itemize}
    \item \textbf{Segmentación Visual (SSIM):} La optimización sobre miniaturas de 160 px alcanza un rendimiento de procesamiento de aproximadamente $45\text{ FPS}$ en CPU Apple Silicon, procesando un vídeo de 4 minutos en menos de $6\text{ segundos}$.
    \item \textbf{Transcripción Whisper Local:} El modelo \texttt{base} transcribe y alinea temporalmente los 4 minutos de audio en aproximadamente $12\text{ segundos}$ en CPU local, operando en paralelo con la segmentación visual.
    \item \textbf{Vectorización en Gemini API frente a Inferencia Local (Etapa 1):} La llamada remota para vectorizar fragmentos con \textit{Gemini Embedding 2} presenta una mediana por consulta de $464\text{--}481\text{ ms}$ (unos $650\text{ ms}$ de media en invocación aislada multimodal y $187\text{--}227\text{ s}$ totales para indexar $82$ fragmentos bajo control de reintentos), mientras que la ejecución local en \texttt{bfloat16} sobre Metal MPS reduce la mediana de consulta a $91\text{--}98\text{ ms}$ en \texttt{qwen3-vl-2b} ($126\text{--}142\text{ s}$ de indexación) y a $\sim 30\text{ ms}$ en \texttt{gemma2} ($2,55\text{--}2,59\text{ GB}$ de memoria Metal), tal como se profundiza en la Subsección~\ref{subsec:comparativa_backends_locales}.
    \item \textbf{Búsqueda Vectorial en Qdrant (Etapa 1):} La consulta sobre el grafo HNSW se completa en $1,14\text{ ms}$ para $d=3072$ y se reduce a \textbf{$0,63\text{ ms}$ para $d=768$} (aceleración del $44,4\%$).
    \item \textbf{Verificación con Árbitro Bimodal (Etapa 2):} La llamada concurrente asíncrona a \textit{Gemini 3.5 Flash Lite} para evaluar y clasificar los 5 mejores candidatos requiere una latencia media de $1,4\text{ s}$, compatible con interfaces de usuario interactivas.
\end{itemize}

\subsubsection{Escalabilidad de la Huella de Almacenamiento Vectorial}
El Cuadro~\ref{tab:escalabilidad_memoria_mrl} proyecta la memoria física requerida para indexar colecciones audiovisuales de diferentes órdenes de magnitud bajo representación en coma flotante de 32 bits (\texttt{float32}).

\begin{table}[htbp]
\centering
\small
\setlength{\tabcolsep}{4.0pt}
\caption{Proyección de memoria vectorial en función del volumen de metraje y la dimensión MRL.}
\label{tab:escalabilidad_memoria_mrl}
\begin{adjustbox}{max width=\linewidth}
\begin{tabular}{lccccc}
\toprule
\textbf{Escala del Catálogo} & \textbf{Fragmentos} & \begin{tabular}[c]{@{}c@{}}\textbf{Duración}\\\textbf{Aprox.}\end{tabular} & \begin{tabular}[c]{@{}c@{}}\textbf{Memoria}\\\textbf{($d=3072$)}\end{tabular} & \begin{tabular}[c]{@{}c@{}}\textbf{Memoria}\\\textbf{($d=1536$)}\end{tabular} & \begin{tabular}[c]{@{}c@{}}\textbf{Memoria}\\\textbf{($d=768$)}\end{tabular} \\
\midrule
\textbf{Banco de Pruebas Actual} & 164 & $\sim 8,1\text{ min}$ & $1,92\text{ MB}$ & $0,96\text{ MB}$ & $\mathbf{0,48\text{ MB}}$ ($-75\%$) \\
\textbf{Curso o Documental} & $10.000$ & $\sim 8,3\text{ horas}$ & $117,19\text{ MB}$ & $58,59\text{ MB}$ & $\mathbf{29,30\text{ MB}}$ ($-75\%$) \\
\textbf{Archivo Departamental} & $100.000$ & $\sim 83,3\text{ horas}$ & $1,14\text{ GB}$ & $585,94\text{ MB}$ & $\mathbf{292,97\text{ MB}}$ ($-75\%$) \\
\textbf{Plataforma Corporativa} & $1.000.000$ & $\sim 833,3\text{ horas}$ & $11,44\text{ GB}$ & $5,72\text{ GB}$ & $\mathbf{2,86\text{ GB}}$ ($-75\%$) \\
\bottomrule
\end{tabular}
\end{adjustbox}
\end{table}

La compresión a $d=768$ dimensiones permite alojar un millón de fragmentos de vídeo (más de 800 horas de contenido) en apenas $2,86\text{ GB}$ de memoria RAM, haciendo viable el despliegue de motores de búsqueda de vídeo de escala empresarial sobre instancias de servidor estándar de bajo coste.

\subsubsection{Discusión y Recomendaciones de Diseño}
La combinación de las evidencias experimentales obtenidas a lo largo de este estudio permite formular una recomendación arquitectónica definitiva:
\begin{enumerate}
    \item En despliegues Cloud (\textit{Gemini}) o \textit{open-weights} de $2\text{B}$ (\textit{Qwen3-VL-2B}), la configuración óptima de compromiso es \textbf{$d = 768$ dimensiones}: retiene el $98,8\%$ de la precisión en primera posición en Gemini (y el $100,0\%$ en Qwen3-VL-2B), mejora el \textit{Recall@5} hasta el $99,02\%$, reduce la huella de memoria un $75\%$ frente a $3072\text{d}$ y acelera las consultas en Qdrant un $44,4\%$. En dispositivos de borde con presupuesto estricto de memoria, \textit{EmbeddingGemma 2} ($512\text{d}$ o $256\text{d}$) aporta una alternativa local de $\sim 30\text{ ms}$ sin coste de API (Subsección~\ref{subsec:comparativa_backends_locales}).
    \item La arquitectura de recuperación en dos etapas resuelve la limitación intrínseca de los modelos de embeddings: la primera etapa aporta velocidad y alta exhaustividad (\textit{Recall@5} $\ge 97\%$), mientras que la segunda etapa aporta capacidad de razonamiento multimodal, eliminando los falsos positivos y elevando la precisión final por encima del $93\%$.
\end{enumerate}
